\chapter{CONCLUSION}
\label{ch:conclusion}

This project developed a mathematical pipeline from unstructured chief complaint text to an ESI-style acuity estimate, a SATS-aligned colour via \(f_{\mathrm{SATS}}\), and a protocol pathway card at hospital intake. The chatbot accepts a single English free-text submission and returns acuity probabilities, a colour recommendation, and a pathway card aligned with SATS practice and healthcare-facility department structure.

\paragraph{What was demonstrated.}
\begin{enumerate}
\item A BioBERT classifier predicts ESI-style urgency from one English chief complaint. Primary reported metrics come from external evaluation on MIMIC-IV-ED and NHAMCS; Triagegeist synthetic pairs are used for training and validation only.
\item A deterministic heuristic \(f_{\mathrm{SATS}}\) maps ESI-style levels to SATS-aligned colours; pathway cards \(P(C_{\mathrm{NLP}})\) name destination, \(T_{\max}\), actions, and escalation, while TEWS remains the nurse vital-sign path in parallel.
\item A clinical gate with threshold \(0.30\) withholds colour from non-clinical chat before BioBERT runs, with Acc/Sens/Spec reported on a labeled contrast set (Section~\ref{sec:gate}).
\end{enumerate}

\paragraph{What the experiments showed.}
External evaluation indicates that free-text chief complaints carry a measurable but limited acuity signal under geographic domain shift. Training with additional Level~1 (Red) examples to offset class imbalance raises Red recall at a clear cost to overall accuracy and precision (RQ4). Softmax mode probability is a model score, not calibrated clinical confidence.

\paragraph{What was not demonstrated.}
This project did not demonstrate live Ghanaian clinical safety, prospective waiting-time reduction, patient-outcome improvement, or clinician-adjudicated SATS equivalence of \(f_{\mathrm{SATS}}\). Pathway tables reflect the authors' documented interpretation of observed workflows and published SATS rules; they are not an official hospital protocol unless formally approved.

Ultimately, this project shows that NLP can be mathematically constrained and operationalized for low-resource hospital settings as nurse-overridable decision support. By bridging unstructured patient narratives to deterministic SATS-aligned pathway protocols at \examplehospital{}, the Flow Management Engine provides an auditable specification with the potential to reduce triage bottlenecks and to support earlier identification of potentially urgent presentations from the moment a patient arrives.

\section{Limitations and future work}

\textbf{Data scope.} Triagegeist is used for Stage~3 training and validation only. Published classification metrics in Chapter~\ref{ch:results} are computed on external MIMIC-IV-ED and NHAMCS text, after dropping empty complaints and retaining acuity labels in \(\{1,\ldots,5\}\), with tokenisation truncated to project length \(M=128\). The model is English-only at intake. An observational site visit to \examplehospital{} mapped pathways and department structure through intake observation and staff inquiries; supplementary English example complaints generated from those protocols support prototype and qualitative analysis only and were not used for fine-tuning. A natural next step is prospective collection of locally labeled chief complaints at \examplehospital{} for local retraining, and exploring optional fusion of vitals and TEWS alongside free text when measurements become available.

\textbf{Label mismatch.} Training and external gold labels are ESI-style; target pathway colours are SATS-aligned via \(f_{\mathrm{SATS}}\). These are related but not identical clinical systems.

\textbf{Cross-domain and cross-linguistic shift.} Formally, \(P_{\mathrm{train}}(X,Y)\neq P_{\mathrm{Ghana}}(X,Y)\). Formal English training text differs from real Ghanaian patient phrasing (Pidgin, Twi-influenced English).

\textbf{Geographic and demographic domain shift.} The chatbot is designed for Ghanaian facility pathways at \examplehospital{}, yet the external test corpora are United States sources. MIMIC-IV-ED is drawn from a Boston-area teaching hospital population that is predominantly White (approximately 68\%), with smaller shares of Black (approximately 15\%), Hispanic/Latino (approximately 5\%), and Asian (approximately 3\%) patients in published MIMIC demographic descriptions. NHAMCS likewise reflects a North American ambulatory and emergency survey population, not Ghanaian intake language mix, care pathways, or case-mix. External accuracy therefore does \emph{not} equal expected live performance at \examplehospital{}; prospective local validation remains essential.

\textbf{Class imbalance and Level~1 scarcity.} Real emergency departments are heavily skewed toward mid-acuity visits. NHAMCS in particular has very few Level~1 (Resuscitation) cases relative to Levels~3--4. Chapter~\ref{ch:results} reports class support \(N_\ell\) alongside recall so rare-class metrics are interpreted cautiously. Stratified oversampling on the Triagegeist training split studies the recall--precision trade-off; classical embedding-space SMOTE was not used for the published oversampled checkpoint (Section~\ref{sec:oversampling-choice}).

\textbf{Age and pediatrics.} The classifier is age-agnostic: it maps \(X\to\hat{y}\), not \((X,\mathrm{age})\to\hat{y}\), and defaults to adult SATS (ASTS) TEWS and pathway tables. Pediatrics is not a default Yellow destination; routing to Pediatrics is reserved for complaints with explicit child markers. Structured age input is future work.

\textbf{No clinical adjudication of colour.} No clinician independently validated chatbot colours against SATS for each external test case.

\textbf{No live workflow or outcome experiment.} Waiting-time reduction and adverse-event rates were not measured.

\textbf{No Softmax calibration study.} Maximum Softmax probability is not established as calibrated clinical probability.

\textbf{Clinical validation.} External offline evaluation suggests text carries a limited acuity signal under domain shift, but Softmax confidence should still trigger nurse oversight when \(\max_c\hat{y}_c < \tau_c\) or when language conflicts with vitals. The clinical relevance gate was validated only on a small hand-labeled contrast set (\(N=52\)), not on prospective \examplehospital{} traffic; lexicon-poor clinical wording can still produce false clinical rejections and must be handled by re-phrasing or nurse override (Section~\ref{sec:gate}). The next phase is a prospective pilot at \examplehospital{} to measure waiting times, nurse override rates, and patient outcomes in live intake \citep{Mitchell2023}.

\textbf{Pathway refinement.} The five-scenario pathway map is SATS-aligned and complete for gate reject, Red, Orange, Yellow, and Green routing, including Green diversion from acute queues. The pathway configuration represents the authors' documented interpretation of the observed workflow and should not be interpreted as an official clinical protocol unless formally approved by the hospital. Ongoing collaboration with hospital administrators and clinical leads at \examplehospital{} will keep department routing tables current as service lines evolve \citep{Rominski2014,satsmanual}.

\textbf{Language and equity.} The current model processes formal English. Expanding the system to handle Ghanaian informal English (Pidgin) or integrating voice-to-text capabilities for Twi would significantly broaden accessibility for patients who are less comfortable writing formal English complaints.

\textbf{System integration.} The prototype chatbot API and pathway cards already anchor the chatbot in a deployable pipeline. Deeper integration with Hospital Information Technology and Health Information systems would automate timer alerts and bed assignment, completing the Flow Management Engine vision across Administration, Diagnostics, and all clinical departments listed in Table~\ref{tab:knust-dept-routing}.
